% Edition 19 audit fragment: first essay 8.14--8.18.
% All source pages and historical source files remain immutable.
\head{D19-8.14--8.18：法英回校、原稿保存与数学订正}
\begingroup
\def\PD{P^\vee}
\def\depth{\operatorname{depth}}
\def\codim{\operatorname{codim}}
\source{法文手稿物理第46页下部至第50页末附记；英编稿第32--34页，止于8.18后的线性系讨论之前}

本项处理已有译稿的法英回校，不是把8.14以后的旧译重新计作新取得原文。实际查看了法文第46--50页及英文第32--34页的图像；法文关键删改和边注另作300 dpi裁图核对。本次不包括法文第46页上半8.13的再次全面审定，也不包括第50页末从“还须明确8.18的条件”起的线性系讨论。

原稿有可读正稿、作者替换、作者删划及英编者疑读；本报告与正文将它们分别标明。以下反例、加强假设版本和补证为本轮译校文字，不归于原作者。旧注原文在本章末逐项保存；旧注中已经由法稿解决的疑读，不再作为本版的未决结论。

\subsec{D19-8.14与8.15：实际读法和保留边界}
\begin{enumerate}
\item 法46页的承接句以手写补入8.6。8.14的开头明确写“$f:X\to P$ 非分歧，$X$ 无孤立点”；下一行的孤立字母 $n$ 不形成条件。英编者第49注“询问A.G.”保留，但不再把这一个残字写成未经核实的数学缺设。
\item 法46页(a)在“充要”后的一段旧条件被框划。可辨的旧层谈维数为1的几何奇异点以及一维的“非直纹”分支；行间添加、连接词和末尾不全，不能将其拼成另一条已采用条件。法47页首另给现行的一维分支条件：由光滑点组成且相对于 $f$ 为“直纹的”。
\item 法47页(a)的深度式只有一次 $\operatorname{prof}_xX\ge2$。英稿重复同一因子，不是原稿在要求一个乘积不等式。法稿的 \emph{séparable} 与本篇3.3一致，读作几何约化；不能再将它列为完全未知、也不能改译成光滑。
\item 法47页(b)的分支维数为 $\le2$，末尾深度为 $\ge3$；其下约三行经密集覆打删除，完整字词仍未恢复。未辨的删除文字保留这种身份，未补作正文。
\item 法47页8.15末明确是 $(R_k)$，英稿的 $E_k$ 是误读。原作者关于 $(S_1)$ 和可能还需等维性的保留意见全部保留。该页正则截面条件后余维数字处于扫描右缘外，现有图像不能恢复这个数值。
\item 法48页 \emph{Remords} 的第一项是在光滑点处非分歧，第二项对充分性只要求在补余维 $\ge k+1$ 的开集上非分歧；这两项不能合并为一个无条件有效的替换规则。其上另有被覆打的旧“命题8.16”起句，可辨 $X$、$P$ 等字母，未恢复完整假设，不作第二个命题输入。
\end{enumerate}

8.14正文继续保留可辨修订层中的充要措辞。本轮确认了其读法，并核对用于8.16加强版本的充分方向；没有宣称重新证明任意非等维情形下的全部必要方向。8.13、8.5和8.6的几何换底、关联点及有限型前提仍须同时使用。尤其不能因8.13已补有伴随点条件，就假称8.14的所有依赖已经自动验证。

\paragraph{原作者补记的完整内容。}
在8.13及8.14中，只须假设 $f:X\to P$ 在 $X$ 的光滑点处非分歧；对于“充分”的方向，更只须它在某个开集 $X'$ 上非分歧，且其补集的余维至少为 $k+1$。

这一段保留为作者草稿提议。以下反例说明，不能仅以所写补集余维界代替相应结论所需的条件。

\subsec{D19-8.16：原陈述和全部作者证明}
\paragraph{原陈述（所写弱假设不足，译校不作为定理采用）。}
假设 $f:X\to P$ 在某个开集上非分歧，该开集的补集余维至少为2；$X$ 几何正规，且在其闭点处深度至少为3；最后，$X$ 几何整且在 $k$ 上固有。那么，使得 $Y_\xi$ 几何正规、几何整且维数为 $\dim X-1$ 的 $\xi\in\PD$ 所成集合（是可构造的，并且）其补集的余维至少为2。

\paragraph{原证明译文。}
事实上，由8.14(b)，对于性质“$Y_\xi$ 几何正规且维数为 $\dim X-1$”，上述结论成立（这一维数性质表达的是 $\varphi_\xi$ 为 $X_{k(\xi)}$-正则）。而由6.1，所有这些 $Y_\xi$ 都几何连通。若 $Y_\xi$ 几何正规，它几何整当且仅当它几何连通，于是完成证明。

法48页在结论前划去了正则截面的重复短语，随后用维数条件表达这一点；原证明的论证长度与最后的连通性步骤均完整保留在上文。问题在于把8.14的非分歧前提弱化后，第一步并未得到有效论证。

\subsec{D19-8.16：余维二收缩反例（本轮译校）}
取 $k=\mathbb C$。在 $\mathbb P^1_{[s:t]}\times\mathbb P^4_{[z:a:b:c:d]}$ 中令
\[
 X=\{at=bs,\ ct=ds\}.
\]
在 $s\ne0$ 的图上，令 $\lambda=t/s$，则 $b=\lambda a$、$d=\lambda c$，故此图为
\[
 \mathbb A^1_\lambda\times\mathbb P^2_{[z:a:c]}.
\]
$t\ne0$ 的图同理。因此 $X\to\mathbb P^1$ 是射影平面丛，$X$ 光滑、整、射影，维数为3，所有闭点的深度为3，特别满足原8.16要求的几何正规、几何整和固有条件。

忘去 $[s:t]$ 得到态射
\[
 \pi:X\longrightarrow Q=\{ad-bc=0\}\subset\mathbb P^4.
\]
$Q$ 的顶点为 $v=[1:0:0:0:0]$，其逆像
\[
 C=\mathbb P^1\times\{v\}
\]
在 $X$ 中余维为2。只要 $(a,b,c,d)\ne0$，秩一矩阵的行方向就唯一恢复 $[s:t]$，并且在某个非零坐标的图上给出正则逆态射；所以 $\pi$ 在 $X\setminus C$ 上是到 $Q\setminus\{v\}$ 的同构。

记 $(x_1,x_2,x_3,x_4)=(a,b,c,d)$。以下31个三次齐次式定义态射 $h:Q\to\mathbb P^{30}$：
\[
 z^3;\qquad
 z x_i x_j\quad(1\le i\le j\le4);\qquad
 x_i x_j x_\ell\quad(1\le i\le j\le\ell\le4).
\]
数量确为 $1+\binom52+\binom63=1+10+20=31$；这里即使在 $Q$ 上存在关系，也不妨碍用这些截面定义到指定射影空间的态射。它们没有公共零点：在 $v$ 处 $z^3\ne0$，其余点至少有一个 $x_i^3\ne0$。

在目标的 $x_i^3\ne0$ 图中，商
\[
 \frac{x_i^2x_j}{x_i^3}=\frac{x_j}{x_i},\qquad
 \frac{z x_i^2}{x_i^3}=\frac z{x_i}
\]
恢复 $Q$ 的原仿射坐标。目标坐标环到 $Q\cap D_+(x_i)$ 的坐标环因而满射，说明 $h$ 在 $Q\setminus\{v\}$ 上为浸入。于是
\[
 f=h\pi:X\longrightarrow\mathbb P^{30}
\]
在 $X\setminus C$ 上非分歧；这恰满足原8.16的开集补余维条件。

然而 $f(C)=[1:0:\cdots:0]$。每个经过此点的目标超平面，其拉回方程都没有 $z^3$ 项。在 $z=1$、$s\ne0$ 的图中写 $a=u,c=w$，便有
\[
 (a,b,c,d)=(u,\lambda u,w,\lambda w),\qquad
 I_C=(u,w).
\]
其余30个坐标全都属于 $I_C^2$。若该拉回方程 $F$ 非零，则在 $C$ 的泛点 $\eta_C$ 处，$A=\mathcal O_{X,\eta_C}$ 是二维正则局部环，而
\[
 0\ne F\in\mathfrak m_A^2,
 \qquad \dim A/(F)=1,
 \qquad
 \dim_{k(\eta_C)}\frac{\mathfrak m_A}{\mathfrak m_A^2+(F)}=2.
\]
故 $A/(F)$ 为一维非正则局部环。一个一维 Noether 正规局部环必为正则局部环，因此该截面不正规；若它非约化，同样不可能正规。若 $F$ 恒为零，则截面等于 $X$，维数为3而不是要求的2。

因此，经过 $f(C)$ 的所有目标超平面都属于例外参数集；这些参数组成 $(\mathbb P^{30})^\vee$ 的一个超平面，余维为1。原8.16的全部假设均满足，结论却失败。本反例同时解释为何只知道非分歧开集的补集余维为2，不足以控制截面在该补集上的正则性。

\paragraph{正文采用的加强版本。}
正文将 $f$ 的假设加强为全局非分歧，其余假设全部保留。深度至少3结合8.6保证截面的 $(S_2)$；几何整且维数至少3结合8.13(b)保证几何 $(R_1)$；前者不能由几何正规本身代替。两类例外的闭包均有余维至少2，有限并仍然如此。由 Serre 判据与6.1，即得正文结论。此处是明确附加条件的充分版本，不称它与原弱假设等价。

\subsec{D19-8.17：恢复的文字与宽泛断言的反例}
\paragraph{原(a)的完整内容。}
8.16的假设蕴涵 $\dim X\ge3$。也许只要 $X$ 几何不可约且 $\dim f(X)\ge3$，不要求正规性和非分歧性，使 $Y_\xi$ 几何不可约的参数集合之补集就具有至少为2的余维。无论如何，可以证明这个补集不包含任何超平面（见下文）。

\paragraph{为何正文不能采用这个宽泛断言。}
取 $f:X=\operatorname{Bl}_p(\mathbb P^3)\to\mathbb P^3$，基域仍取 $\mathbb C$。$X$ 光滑、整、射影，$\dim f(X)=3$。设 $E\simeq\mathbb P^2$ 为例外除子。若超平面 $H$ 经过 $p$，其局部方程在 $p$ 处为一阶消失，故
\[
 f^*H=\widetilde H+E.
\]
$\widetilde H\simeq\operatorname{Bl}_p(\mathbb P^2)$ 与 $E$ 是不同的非空不可约除子，因此总截面可约。过 $p$ 的超平面参数恰构成 $(\mathbb P^3)^\vee$ 中的一个超平面，推翻了原文的“无论如何”断言，也推翻其所猜测的余维至少2结论。这个例子的非分歧例外是一个除子，不是有限点集，故不能把它误称为8.18的反例。

\paragraph{二次曲面例子。}
法48页明确有“相切超平面”；英稿缺文已经恢复。以特征零光滑二次曲面为例，在代数闭域上可写 $x_0x_3-x_1x_2=0$。在点 $[1:0:0:0]$ 处的切平面为 $x_3=0$，截面方程为 $x_1x_2=0$，即两条相交直线。切平面参数构成对偶二次曲面，维数为2、余维为1。本计算仅核实这个反例，不把法稿关于“对偶形式”的短句提升为任意特征下的全部对偶理论。

\paragraph{法49页顶端的手写(c)。}
此段可辨的完整条件为：$X$ 光滑、二维、几何整、固有，$f$ 为浸入；引用6.1、8.8，旁补8.14；结论为可分截面至多有两个几何不可约分支，例外余维至少2。关于 $g:Y^{\mathrm{sing}}\to\PD$ 的原条件写作 \emph{n'est pas gén. fini et inséparable}。“不”的作用范围没有明确另列逻辑式；本版既不改作已确定的“非泛有限或可分”，也不改作“非泛有限且不可分”。此外8.8起首的缺文仍在。正文把它标为作者未定稿的手写补充，不宣称完成了该一般结论的独立证明。

\subsec{D19-8.18：原完整陈述与实际法语读法}
\paragraph{原陈述（按所印条件不能采用）。}
假设 $k$ 代数闭，$X$ 整（分别地，整且正规），维数 $d\ge2$，且在 $k$ 上固有。设 $T$ 为 $X$ 的有限闭子集，使 $X-T$ 光滑，并使 $f|_{X-T}$ 非分歧。为了使满足 $Y_\xi$ 几何不可约（分别地，几何整）且“维数为1”的 $\xi\in\PD$ 所成集合，其补集余维至少为2，原稿称以下条件为充要条件：
\begin{enumerate}
\item[(a)] 对每个 $x\in T$，存在一个经过 $x$、维数为 $d-1$ 且不可约的超平面截面 $Y_\xi$，$\xi\in\PD(k)$。
\item[(b)] $X'=X-T$ 相对于 $f$ 是“直纹的”（原稿如此），或者存在 $X'$ 的一个维数为 $d-1$、奇异且不可约的超平面截面 $Y'_{\xi_0}$，$\xi_0\in\PD(k)$。
\end{enumerate}

法49页首句仍印“维数1”；(a)及法50页证明中的相应维数已由作者手改为 $d-1$。因而这不只是英编稿的打字错。若按任意 $d\ge2$ 的原假设读首句，取 $X=\mathbb P^3$、$f=\mathrm{id}$、$T=\varnothing$，所有超平面截面维数为2；首句维数1的结论直接失败。正文采用 $d-1$，明确标作译校采用版本，不冒称原首句已如此改定。

法49页(b)最后一行清楚为“奇异且不可约”；前行有一处被覆打的旧词，但末行 \emph{singulière} 未被否定。英编者第51注的“(non???)”疑读可以解决，不能继续把“非奇异”列作同等可用读法。(a)参数误作 $P(k)$，由上下文及(b)统一为 $\PD(k)$。

\paragraph{原稿全部证明译文（保留其草稿身份）。}
先假设 $X$ 几何正规。由8.14(a)，我们已经知道可以在 $\PD$ 中找到一个余维至少为2的闭子集 $Z'$，使得对 $\xi\in\PD-Z'$，$Y_\xi$ 在 $k(\xi)$ 上可分且维数为 $d-1$。对这样的 $\xi$，$Y_\xi$ 几何不可约与几何整等价，所以8.18所考察的两个问题（带“分别地”与不带“分别地”）等价。

另一方面，由5.6，使 $Y_\xi$ 几何整且维数为 $d-1$ 的 $\xi\in\PD$ 所成集合 $U$ 是开的（这里维数假设表达的是 $\varphi_\xi$ 为 $X_{k(\xi)}$-正则）。

我们将给出一个包含于 $U$ 的、明显的非空开子集 $\PD-Z$：取 $Z$ 为
\[
 \overline{g(Y'^{\mathrm{sing}})}
 \quad\text{与各 }H_x\ (x\in T)\text{ 的并},
\]
其中 $H_x$ 是由 $f(x)$ 确定的对偶参数超平面。对于 $\xi\in\PD-Z$，$Y_\xi$ 光滑且维数为 $d-1$；又由6.1，它几何连通，因而几何整。

因此，我们要把 $Z$ 的每个余维一不可约分支都与开集 $U$ 相交这一条件表达出来。这些不可约分支就是各 $H_x$（可能重复，但无关紧要），以及当 $\overline{g(Y'^{\mathrm{sing}})}$ 的余维确实为1时的这个集合；后者对应于 $X'$ 相对于 $f$“不是直纹的”的情形。（注意：这里使用了 $Y'^{\mathrm{sing}}$ 的不可约性。）

再者，为使后一个集合与开集 $U$ 相交，充要条件是 $g(Y'^{\mathrm{sing}})$ 与 $U$ 相交，因为它在其闭包中包含一个开稠密集。这便证明了这一情形下的8.18。如果不假设 $X$ 正规，就对 $X$ 的正规化应用前面的结果；论证是直接的，我在这里不写细节。

注意：在带“分别地”的情形，8.18包含在8.16中，更确切地说，除 $d=2$ 以外如此。对不带“分别地”的情形，也许有必要不要求 $d=2$，而只要求 $d\ge2$。

\paragraph{证明中的原字与缺口。}
法50页前两次出现的横线都清楚盖在整个像 $g(Y'^{\mathrm{sing}})$ 上方；最后一次是没有上横线的像本身，随后括号说它在闭包中含开稠密集。旧英文排版及旧中文注释对此判断失准，已由本次300 dpi裁图解决。$Y'$ 下方的短横不应误作闭包的上横线。正文采用的像闭包记号有原图依据，不是凭拓扑常识补写。

原证明的真正数学缺口是：原(a)、(b)只给出不可约截面，不保证它们对应的参数属于几何整的开集 $U$。在一个坏除子内找到特殊的非约化但不可约截面，并不能排除这个除子的泛点仍然可约。下面的两个反例分别检验这一点。原末尾的正规化归约没有写出条件如何传递；末尾附记也没有补入8.16的闭点深度假设。本版不把这两句作者省略声明当作已经完成的证明。

\paragraph{被划去的证明层。}
法49页下部整段框划的旧证明，可辨内容是：令 $Z$ 为切关联像和各 $H_x$ 的并；对补集参数，截面光滑且维数为1；再由6.1的几何连通性得到几何不可约。中间两行有连续覆打，不能可靠恢复。它没有第50页新证明中的余维一分支逐一相交论证；不能把这块删稿与新证明合成一段作者定稿。

法50页首在“先假设 $X$ 几何正规”之前还有一条被密集覆打删除的旧行，可辨出维数及几何整性的字样，但完整命题、关系符号和连接词未恢复。第49页(b)前行的旧词及这两处覆打，作为真正未辨处继续登记，未在正文中补造。

\subsec{D19-8.18：原条件的反例与充分版本}
\paragraph{反例一：二次 Veronese 曲面检验(b)。}
取 $k=\mathbb C$、$X=\mathbb P^2$，令 $f=\nu_2:X\hookrightarrow\mathbb P^5$ 为全部二次单项式定义的 Veronese 嵌入，取 $T=\varnothing$。$X$ 光滑、正规、整、射影，且 $f$ 为闭浸入，满足原稿全部环境假设。(a)为空条件；取二次曲线 $x^2=0$，便得到维数1、奇异且拓扑不可约的截面，所以原(b)成立。

然而超平面参数为平面二次曲线的系数。以
\[
 q=a x^2+b y^2+c z^2+2dxy+2exz+2fyz
\]
记二次式，其判别式为
\[
 \Delta=\det\begin{pmatrix}a&d&e\\d&b&f\\e&f&c\end{pmatrix}
 =abc+2def-af^2-be^2-cd^2.
\]
$D=V(\Delta)\subset\mathbb P^5$ 是所有奇异二次曲线的不可约闭集：它由两条直线的无序对参数化，故维数为4。一般点对应两条不同直线之并，因而几何可约。二倍直线只形成一个维数2的闭子集，不能排除 $D$ 的泛点为坏参数。

所以，几何整截面的坏集合为 $D$；即使只要求几何不可约，坏集合仍包含 $D$ 中稠密开集，闭包也为 $D$。两种结论的坏参数余维都为1，而不是2。这个反例在 $d=2$ 时已经成立，不依赖首句“维数1”的另一个错误。

\paragraph{反例二：正规二次锥面检验(a)。}
在 $\mathbb P^3_{[x:y:z:w]}$ 中令 $X=V(xz-y^2)$，$p=[0:0:0:1]$ 为顶点，$f$ 为闭浸入，$T=\{p\}$。这是正规整射影曲面：它是 Cohen--Macaulay 超曲面，唯一奇异点余维为2，故满足 $(S_2)$ 与 $(R_1)$。

原(a)可取 $x=0$ 的截面；此时 $y^2=0$，截面拓扑支撑是一条直线，故不可约且维数1，却非约化。$X'=X\setminus\{p\}$ 的切平面参数为
\[
 [z:-2y:x:0],\qquad xz-y^2=0,
\]
其像闭包是对偶空间中的一条圆锥曲线，维数1、余维2，因而满足原(b)的“直纹”分支。

但一般经过 $p$ 的超平面与锥面的交是两条不同母线之并；过 $p$ 的参数超平面维数2。特殊的切平面给出二倍母线，仍不能使几何整参数在这个除子中稠密。因此原(a)的存在性同样不能保证 $H_p\cap U\ne\varnothing$。

\paragraph{正文充分版本的界限。}
正文只处理正规 $X$，把(a)中的整体截面、(b)中的开截面均加强为几何整，并明确第一分支为 $\codim\overline{g(Y'^{\mathrm{sing}})}\ge2$。它只宣称充分，不宣称恢复了原作者的一般充要命题。

关键补足有两项。首先，固有平坦有限表示族的几何整纤维参数是开集；不能以几何不可约参数的一般可构造性替代这个开性。几何约化纤维开性的现代出处可核为 \href{https://stacks.math.columbia.edu/tag/0C0E}{Stacks Project 0C0E}；几何不可约性在一般有限表示族中只有可构造性，须与约化及此处的固有平坦条件一起使用。其次，正规环境中的有效 Cartier 除子为 $(S_1)$，所有分支均有维数 $d-1$。删去有限集 $T$ 后若开截面几何整，则原截面没有在 $T$ 中的额外分支，也无嵌入伴随点，因而仍几何整。

余维一的闭包 $D$ 由不可约切关联概形支配；条件(b)的奇异几何整截面确实给出 $D\cap U$ 的点，故 $D$ 的泛点也在开集 $U$ 中。各 $H_x$ 同理由(a)处理。这正补上原证明误把“不可约”当成足以落入 $U$ 的那一步，而不采用它未详述的正规化推广。

\subsec{本轮实际检查与仍未完成的事项}
\begin{itemize}
\item 实际图像范围：法46--50及英32--34；关键裁图为法46页8.14、法47页8.15末、法49页(c)与8.18(b)、法50页闭包及末附记。
\item 实际读法恢复：8.14开头、深度式、8.15的 $(R_k)$、8.17相切截面和手写(c)、8.18的“奇异”、三处像闭包写法、末附记的“只要求”。
\item 数学核验：8.16的余维二反例、8.17(a)的吹起反例、8.18的 Veronese 和锥面反例；正文两个加强版本按所列条件给出独立补证。代数单项式与判别式的实际符号核算记录另存本目录JSON。
\item 保留未决：8.14任意非等维情形下全部必要方向未在本次重证；8.15右缘缺失数值；若干密集覆打删除层；8.17(c)否定的逻辑作用范围及8.8缺失起首；8.18不要求正规的原一般版本及原末尾包含关系未建立。
\item 本子任务不构建整本PDF。模块编译、单文件展开和最终版面由总工程构建记录决定，不以本报告宣称这些后续步骤已经通过。
\end{itemize}

% BEGIN D19 HISTORICAL NOTES
\subsec{D19-8.14--8.18：第十八版旧注逐条原样留存}
下列各项按第十八版原文件中的出现顺序保存旧脚注、译注和编者注，注内文字与数学源码逐字保留，仅把脚注容器改为本报告的编号条目。它们是历史记录，不是第十九版新的待核清单：其中关于“相切”、$(R_k)$、8.17(c)的可辨条件、8.18的奇异性、闭包横线及末附记的旧疑读，已经由本轮法文图像核对更新。应以本报告前述校记及正文现行脚注为准。原片段另存本目录的文本档案，可逐字复核正文与注释的旧状态。
\begin{enumerate}
\item \textbf{旧译注 1。}此为原作者的引导句；底本未另列完整推导，本译不把它当作已经进行独立验证的声明。
\item \textbf{旧脚注 2。}原脚注49：“询问 A.G.”。
\item \textbf{旧译注 3。}假设末尾的残字 $n$ 及缺文照录，无法确定其数学含义，不据此补造附加假设。
\item \textbf{旧译注 4。}“直纹的”是原稿的暂用术语，不能径同于通常直纹簇概念；底本8.17专门指出二者不一致。本条沿用第十七版译稿，完整保存原陈述的充要方向。第十八版已用第108批回校稿替换8.13，并明示回校采用的新增假设；这项替换不等于本条8.14及其依赖关系已经重新审定，故这里仍不提供未经核实的补证。
\item \textbf{旧译注 5。}底本用词为 \emph{separable}。此处按原用语译为“可分”，不擅自改成“光滑”或“几何正规”；该旧用语在全书中的精确定义仍须结合前文核定。
\item \textbf{旧译注 6。}底本将 $\phi_\xi$ 误排成 $\phi_x i$，按本节一贯记号订正。深度条件印作 $\operatorname{prof}_x X\ge2\operatorname{prof}_x X$，重复了式末的 $\operatorname{prof}_x X$；正文删除该重复。依据是本句随后关于几何正规概形的说明，以及(b)的平行深度条件。按字面把它解释成乘积会使正深度点不满足条件，与原文紧随的说明不相容。此为排印校订，不是宣称已验证整个充要条件。$\operatorname{prof}$ 表示深度。
\item \textbf{旧译注 7。}这一保留意见为原文的一部分，必须与前面“不限制一般性”的断言一并阅读。底本的 $E_k$ 记号及缺失的余维数字均未恢复；不能把本段改写成无条件成立的去假设定理。
\item \textbf{旧译注 8。}底本将这段另题为“Remorse”，这里译作作者重新考虑前述假设的补记。它没有展开验证，且 $k$ 在8.14不同性质中的取值未在此重述；本译保留原范围和充分性限定，不据此全面替换前面条目的假设。
\item \textbf{旧译注 9。}8.14前后的假设调整见8.15后的作者补记；其原文缺失和依赖问题仍未全部恢复。此处保留作者引用，不把引用本身作为此处译校已独立验证的证据。
\item \textbf{旧译注 10。}底本在此夹有法文“dès que”的转写痕迹。正文译为“只要”，并保留作者“It is possible”的推测语气，不将其升级为定理。
\item \textbf{旧译注 11。}原句代词“it”依上下文指例外补集；此处明确所指，不改变其断言范围。
\item \textbf{旧脚注 12。}原脚注50：“询问 A.G.”。
\item \textbf{旧译注 13。}底本此处印作“在 $P$ 中”，但所计数的是超平面的参数族，且紧随文字明确提到对偶空间，故订为 $P^\vee$。缺文处未补入“相切”等字样；这里的“这些超平面截面”只承接原稿所讨论的可约截面，不是说该曲面的所有超平面截面都可约。
\item \textbf{旧译注 14。}本段从假设到结论均有缺字和断句。“not generically finite and inseparable”的否定范围亦不明确；正文保留其不确定性，不自行判定应否定合取，或仅否定“泛有限”。不能由这些残句恢复一个确定的充分条件或结论。8.8亦尚缺完整假设，故不补造推论。“separable”沿用前批译为“可分”。
\item \textbf{旧译注 15。}“显然”是原作者的表述，底本此处未给证明；此处译校不另加未经核对的论证。
\item \textbf{旧译注 16。}这是作者对暂用术语的自我修正，不是关于通常意义下二次曲面是否直纹的否定。
\item \textbf{旧译注 17。}底本此处确实写“dimension 1”，而(a)、(b)和随后的证明均用 $d-1$。末尾附记又讨论是否应限于 $d=2$，故不能仅凭上下文断定是漏写 $d-$ 还是原拟限定曲面。正文明确引述这一未决读法，不把它作为已校定的普遍定理。这里设 $d=\dim X$，只是明示原稿随后使用的记号。
\item \textbf{旧译注 18。}底本(a)印作 $P(k)$，但 $\xi$ 参数化超平面，命题首句及(b)均使用对偶空间；这里统一为 $P^\vee(k)$。下文作为参数空间出现的 $P$ 同样统一为 $P^\vee$。
\item \textbf{旧脚注 19。}原脚注51：“字迹不清，询问 A.G.”。
\item \textbf{旧译注 20。}底本写“(non???) singular”。究竟要求奇异还是非奇异，现有文本不能确定；这里不择其一。“直纹的”是作者暂用术语，应连同8.17后关于该术语的自我修正一起理解。
\item \textbf{旧译注 21。}底本此处像集未标闭包；正文不增添横线。$X\in T$ 的大写 $X$ 订为点变量 $x\in T$，与同句 $f(x)$、$H_x$ 及(a)一致。
\item \textbf{旧译注 22。}底本这第三处横线位于 $Y'^{\mathrm{sing}}$ 上方，与前两处记号不同；不能未经原稿核对将其改为像的闭包。底本这句括号没有正常闭合，“which contains an open and dense set”的指代不够清楚。正文保留其内容而不擅自补入一个新的拓扑断言。按通常可构造像的论证，可以考虑像在其闭包中包含开稠密子集，但这里不把这种解释认定为手稿缺文的恢复。
\item \textbf{旧译注 23。}这是作者的省略声明，不是此处译校完成了正规化步骤的全部假设核验；尤其原稿(b)的奇异性条件仍待核。
\item \textbf{旧译注 24。}最后这句话按底本“not to require $d=2$ and not only $d\ge2$”保留，其否定与比较关系在当前转写中不清楚。不能据此确定应增加还是删除曲面限制，亦不能据此前文维数 $1$ 的歧义宣称已修复。
\end{enumerate}
% END D19 HISTORICAL NOTES
\endgroup
